\SourceLeaf{70}
Sei $r=2m-1$. Dann ergibt sich

\[
\frac{1}{2^{r-1}} = \frac{1}{2^{2m-2}},
\]

\noindent und wegen

\[
\frac{1-(-1)^r}{4} = \frac{1-(-1)^{2m-1}}{4} = +\frac{1}{2},
\]

\[
{}^{r}\mathfrak{B}^{\left(\frac{1}{2}r+\frac{1-(-1)^r}{4}\right)} = {}^{2m-1}\mathfrak{B}^{\left(\frac{1}{2}(2m-1)+\frac{1}{2}\right)} = {}^{2m-1}\mathfrak{B}^{(m)},
\]

\noindent schließlich

\[
\cos\frac{1-(-1)^r}{2}\varphi = \cos\frac{1-(-1)^{2m-1}}{2}\varphi = \cos\varphi.
\]

Schließlich ist der vorangestellte Faktor für ungerades $r$ $=1$,
für gerades $=1-\tfrac{2}{4}=\tfrac{1}{2}$, wie es sein muss.

Ebenso behaupte ich, dass für $\cos\varphi^r$ ein beliebiges Glied $h^{\text{tum}}$ lautet

\[
= \frac{1}{2^{r-1}}\,{}^{r}\mathfrak{B}^{\left(\frac{1}{2}r+\frac{1-(-1)^r}{4}+h\right)}\cos\left(2h+\frac{1-(-1)^r}{2}\right)\varphi.
\]

Ist nämlich $r$ gerade, $=2m$, so ergibt sich wegen

\[
\frac{1-(-1)^r}{2} = 0,
\]

\noindent der Ausdruck

\[
\frac{1}{2^{2m-1}}\,{}^{2m}\mathfrak{B}^{(m+h)}\cos2h\varphi;
\]

\noindent ist es ungerade, $=2m-1$, so gilt dann

\[
\frac{1-(-1)^{2m-1}}{2} = 1,
\]

\noindent und derselbe Ausdruck wird zu

\[
\frac{1}{2^{2m-2}}\,{}^{2m-1}\mathfrak{B}^{(m+h)}\cos(2h+1)\varphi,
\]

\noindent wie es sein muss.

Um die Schreibarbeit zu erleichtern, könnte man für die einfache Größe

\[
\frac{1-(-1)^r}{2}
\]

\noindent ein einziges Zeichen, z.\,B. $=a$, verwenden, wie es in der Analysis üblich
ist. In kürzerer Schreibweise erhielte man dann
